%% ASCE-style engineering manuscript starter — ascelike document class.
%%
%% ascelike is NOT included in this download: it ships with every modern TeX
%% distribution (and with latextemplates.com's compiler), so bundling a copy
%% would only risk pinning you to an older version than your own system
%% already has.
%%
%% ascelike is an UNOFFICIAL class that approximates the American Society of
%% Civil Engineers manuscript format. Check ASCE's current author guidelines
%% before submitting.
%%
%% Class documentation: texdoc ascelike   |   https://ctan.org/pkg/ascelike
%% Licence: ascelike is LPPL-1.3, (c) Matthew R. Kuhn.
%%
%% Class options: `Review' gives the double-spaced review layout ASCE asks for
%% at submission; `NewProof' gives the camera-ready single-spaced form.
\documentclass[Review]{ascelike}

\usepackage[T1]{fontenc}
\usepackage{graphicx}

\NameTag{Lovelace, \today}

\begin{document}

\title{Your Paper Title Here}

%% ascelike uses authblk: authors and affiliations are declared separately,
%% with \affil following the \author(s) it applies to. Do NOT use \thanks or
%% \and here — authblk collects the lists itself and those break \maketitle.
\author[1]{Ada Lovelace}
\author[2]{Second Author}
\affil[1]{Analytical Engines Laboratory, Institute of Analytical Engines,
London, United Kingdom. E-mail: ada.lovelace@example.edu}
\affil[2]{Department of Civil Engineering, Another University, Springfield, USA.}

\maketitle

\begin{abstract}
Replace this with your abstract. A good abstract states the problem, what you
did, and what you found --- in that order, and without citations.
\end{abstract}

\KeyWords{First keyword; second keyword; third keyword}

\section{Introduction}
Start here. Cite things like this~\cite{knuth1984texbook}.

ASCE journals expect a review-stage manuscript to be double spaced with
numbered lines, and the camera-ready version to be single spaced. The class
option at the top of this file switches between the two.

\section{Method}
Describe what you did.

\section{Results}

\begin{table}
\caption{A caption belongs above a table.}
\label{tab:results}
\centering
\begin{tabular}{lrr}
\hline
Method & Capacity (kN) & Deflection (mm)\\
\hline
Baseline & 118 & 4.2\\
Proposed & 143 & 3.1\\
\hline
\end{tabular}
\end{table}

Table~\ref{tab:results} reports the comparison.

\section{Conclusion}
Summarise the contribution and name the next step.

\bibliographystyle{ascelike}
\bibliography{references}

\end{document}
